\section{Connectivity at weights growing with the prime}
\label{sec:connectivity}

Fix a stage $0\leq a<n$ and the universal coefficient algebra and
finite-cell module constructed in Section~\ref{sec:universal}.
Equation~\eqref{eq:whole-module} describes $\gr K_p^a$ as a whole
graded module on a set of primes in $\U$. We will use that description
at the moving weight $p\lambda+\delta$, where $\lambda,\delta\in\Lambda$
are fixed. The goal is a lower degree bound with explicit linear
dependence on $p$ after subtracting the ordinary sphere degree
$d(p\lambda+\delta)$.
The resulting vanishing, after tensoring with $K^a$, is the input to
the normalization construction in Section~\ref{sec:normalization}.

We say that a spectrum belongs to $\Sp_{\geq L}$ if its homotopy
groups vanish in degrees strictly below $L$. All spectra in the
primewise calculations below are $p$-local. Sphere degrees may be
negative; the arguments apply to virtual Thom spectra after an
ordinary integer suspension.

\subsection{Symmetric groups and multiplication cofibers}

For an integer $d$, write
\[
 \Sym^k(S^d)=\big((S^d)^{\otimes k}\big)_{h\Sigma_k}.
\]
This is the Thom spectrum of the virtual bundle $d\rho_k$ over
$B\Sigma_k$, where $\rho_k$ is the real permutation representation.
Its virtual rank is $kd$, and its mod-$p$ orientation character is
$\operatorname{sgn}^d$. Thus its mod-$p$ cohomology, with degrees
translated by $kd$, is group cohomology with trivial coefficients
when $d$ is even and with sign coefficients when $d$ is odd.

\begin{lemma}\label{conn:stable-elements}
Let $p$ be odd, let $0\leq k<p^2$, and write $k=lp+b$ with
$0\leq b<p$. Let $H=(C_p)^l\subseteq\Sigma_k$ act by translations
on $l$ disjoint blocks of $p$ letters, and let $N$ be its normalizer.
For either the trivial or the sign module $M$ over $\F_p$, restriction
identifies
\[
 H^*(\Sigma_k;M)
 \xrightarrow{\ \cong\ }
 H^*(H;M)^{N/H}.
\]
Here the normalizer acts both by conjugation on $H$ and by its given
action on $M$. Moreover,
\[
 N/H\cong
 \big((\F_p^\times)^l\rtimes\Sigma_l\big)\times\Sigma_b.
\]
\end{lemma}

\begin{proof}
The exponent of $p$ in $k!$ is $l$, since $k<p^2$.
Consequently $H$ is a Sylow subgroup. Its nontrivial orbits are
exactly the $p$-element blocks, so a normalizing permutation permutes
these blocks and the $b$ remaining letters separately. On each
block the normalizer of translations consists of affine maps
$x\mapsto ux+v$. Dividing out the translations gives the displayed
normalizer quotient. Its order is $(p-1)^l l!b!$, which is prime
to $p$ because $l,b<p$.

We recall the transfer identities needed here, including their
effect on coefficients. Cohomological transfer is obtained by
summing a cochain over representatives of the relevant finite
coset set, using the group action to identify the coefficient
fibers. It follows that transfer after restriction from a group to
a subgroup is multiplication by the subgroup index. For the
opposite composite, decomposing that coset set into subgroup orbits
gives the double-coset formula: each double coset contributes
restriction to the corresponding intersection, conjugation with
its coefficient action, and transfer from that intersection.
These descriptions also give the projection formula for transfer
and cup products.

Since $[\Sigma_k:H]$ is prime to $p$, transfer after restriction
proves that restriction to $H$ is injective. Its image is invariant
under the normalizer: conjugation together with the coefficient
action acts trivially on cohomology of the full group, and
restriction respects this action.

To prove that every invariant is in the image, consider the other
transfer composite. The restriction of either $M$ to $H$ is trivial:
each $p$-cycle is even. If $L$ is a proper subgroup of $H$, it is a
proper linear subspace of the elementary abelian group $H$. A
linear retraction $H\to L$ shows that
$H^*(H;\F_p)\to H^*(L;\F_p)$ is surjective. Choose a basis of the
one-dimensional, trivial $H$-module $M$. Every class with
coefficients in $M$ over $L$ is then the restriction of a class
over $H$. The projection formula gives
\[
 \operatorname{tr}_L^H\operatorname{res}_L^H(x)
   =x\,[H:L]=0.
\]
Hence transfer from every proper intersection subgroup in the
double-coset formula vanishes, in all degrees. An intersection
$H\cap gHg^{-1}$ equals $H$ exactly when $g$ normalizes $H$.
The remaining terms are therefore
\[
 \operatorname{res}_H^{\Sigma_k}
 \operatorname{tr}_H^{\Sigma_k}(x)
   =\sum_{\bar g\in N/H}\bar g\cdot x.
\]
For an invariant $x$ this is $|N/H|x$. The scalar is invertible in
$\F_p$, proving surjectivity onto the invariants.
\end{proof}

For even $d$ define the multiplication cofibers
\begin{equation}\label{conn:cofiber-definition}
 T_{p,0}(d)=S^0,\qquad
 T_{p,k}(d)=\cofib\left(
 S^d\otimes\Sym^{k-1}(S^d)
 \longrightarrow\Sym^k(S^d)\right)\quad(k\geq1).
\end{equation}
The arrow is multiplication in the free commutative algebra on
$S^d$. In the homotopy-orbit description of its symmetric powers,
this multiplication comes from the block inclusion
$\Sigma_1\times\Sigma_{k-1}\subseteq\Sigma_k$. Thus the arrow in
Equation~\eqref{conn:cofiber-definition} is induced by
$B\Sigma_{k-1}\to B\Sigma_k$, with the compatible Thom bundles.
In particular its map on cohomology is restriction. This
identification specifies the map itself and its bottom-degree
normalization; no transfer or factor of $k$ is inserted.

\begin{proposition}\label{conn:symmetric}
Fix a positive integer $C$. For every odd prime $p>C$, every
$0\leq k\leq Cp$, and every integer $d$, one has
\[
 \Sym^k(S^d)\in\Sp_{\geq kd}.
\]
If $d$ is odd, there is the additional bound
\begin{equation}\label{conn:odd-bound}
 \Sym^k(S^d)\in\Sp_{\geq kd+k-(2C+1)}.
\end{equation}
If $d$ is even, the objects in
Equation~\eqref{conn:cofiber-definition} satisfy
\begin{equation}\label{conn:cofiber-bound}
 T_{p,k}(d)\in\Sp_{\geq kd+2k-3C}.
\end{equation}
More precisely, writing $k=lp+b$ with $0\leq b<p$, an odd-degree
extended power is contractible if $b\geq2$, and otherwise its
cohomology above its rank begins no earlier than $l(p-2)$.
For even $d$ and $k>0$, the multiplication cofiber is contractible
if $b>0$, and if $b=0$ its cohomology above its rank begins no
earlier than $l(2p-3)$.
\end{proposition}

\begin{proof}
The Thom spectrum has a cellular filtration with its cells in
degrees $kd+j$, where $j\geq0$ is a base-cell degree. This proves
the first bound, including for virtual bundles. Since $p>C$ and
$k\leq Cp$, we have $k<p^2$ and $l\leq C<p$, so
Lemma~\ref{conn:stable-elements} applies.

The two-periodic cyclic-group resolution, together with its
periodicity class and Bockstein, gives
\[
 H^*(C_p;\F_p)=\Lambda(\eta)\otimes\F_p[\xi],
 \qquad |\eta|=1,\quad |\xi|=2.
\]
Both $\eta$ and $\xi$ transform with multiplier weight one under
$\F_p^\times$, or both with its inverse if the conjugation
convention is reversed. The periodic resolution gives the
polynomial degree-two class, the Bockstein identifies its
multiplier action with that of the degree-one class, and
$\eta^2=0$ follows from graded commutativity since $p$ is odd.

The sign of multiplication by $u\in\F_p^\times$ on the $p$
letters of a block is $u^{(p-1)/2}$, viewed in $\F_p$.
Indeed, applying the permutation to the Vandermonde product on
the $p$ field elements multiplies that product both by its sign
and by $u^{p(p-1)/2}=u^{(p-1)/2}$. Translations have sign one.
A block transposition has sign $(-1)^p=-1$, and permutations of
the remaining $b$ letters have their ordinary sign.

Suppose first that $d$ is odd. If $b\geq2$, a transposition of
the remaining letters acts trivially on $H$ and by $-1$ on the
coefficient line. The normalizer invariants are consequently zero
in every degree. If $b=0$ or $1$, invariance under the independent
multipliers on each block requires multiplier weight
$(p-1)/2$ modulo $p-1$ in each factor of
$H^*(H;\F_p)$. The least degree with this weight is $p-2$,
represented by $\eta\xi^{(p-3)/2}$; the least even degree is
$p-1$. Reversing the multiplier convention gives the same
condition. Therefore the normalizer invariants have no classes
below degree $l(p-2)$. The further block-permutation condition
can only remove classes. Since $b\leq1$ in this case,
\[
 l(p-2)=k-(2l+b)\geq k-(2C+1).
\]
This gives the cohomological vanishing required for
Equation~\eqref{conn:odd-bound}.

For even $d$, put
\[
 U=H^*(C_p;\F_p)^{\F_p^\times},\qquad
 V=\ker(U\longrightarrow\F_p).
\]
The least positive degree in $U$ is $2p-3$, represented by
$\eta\xi^{p-2}$. By Lemma~\ref{conn:stable-elements}, trivial
coefficients give
\[
 H^*(\Sigma_k;\F_p)\cong(U^{\otimes l})^{\Sigma_l}.
\]
The permutations of the tensor factors here include the Koszul
signs of their cohomological degrees. The leftover $\Sigma_b$
acts trivially and imposes no further condition.

We now compute restriction to $\Sigma_{k-1}$. If $b>0$, choose
the fixed letter outside the $p$-blocks. Both symmetric groups
then have the same Sylow subgroup $H$, and their normalizer
invariants are the same: the only change is from $\Sigma_b$ to
$\Sigma_{b-1}$ on the unused letters. The restriction is an
isomorphism, since its further restriction to $H$ is the identity.

If $b=0$ and $k>0$, choose the fixed letter in the last block.
A Sylow subgroup of $\Sigma_{k-1}$ consists of the first $l-1$
blocks. Naturality of restriction and its injectivity on these
Sylow subgroups identify the restriction map with
\begin{equation}\label{conn:restriction}
 (U^{\otimes l})^{\Sigma_l}
 \longrightarrow (U^{\otimes(l-1)})^{\Sigma_{l-1}},
 \qquad \id^{\otimes(l-1)}\otimes\epsilon,
\end{equation}
where $\epsilon:U\to\F_p$ is augmentation.
To determine this map without a choice of normalization, decompose
$U=\F_p\oplus V$. A tensor with exactly $r$ entries from $V$
has its nonunit entries in an $r$-element subset of $\{1,\ldots,l\}$.
An invariant on all such subsets is uniquely determined by its
component on one subset, which belongs to
$(V^{\otimes r})^{\Sigma_r}$. Conversely, place such a component
in each $r$-element subset and sum the resulting tensors, without
dividing the sum by its number of terms. This gives all the
invariants and respects the graded signs. Applying augmentation
to the last factor kills exactly the summands whose subset
contains that factor. For $r<l$, the remaining sum is precisely
the corresponding sum over the $r$-element subsets of
$\{1,\ldots,l-1\}$. For $r=l$, every summand is killed.
Consequently Equation~\eqref{conn:restriction} is surjective and
its kernel is $(V^{\otimes l})^{\Sigma_l}$. Every class in this
kernel has degree at least $l(2p-3)$.

The Thom isomorphism is compatible with the block inclusion
defining $T_{p,k}(d)$: on $\Sigma_{k-1}$ the bundle $d\rho_k$
restricts to $d\rho_{k-1}$ plus a trivial bundle of rank $d$.
Both Thom spectra therefore have rank $kd$. The preceding
restriction maps are surjective in every degree. The cohomology
long exact sequence of their cofiber identifies its cohomology
with the kernel of restriction in the same degree, without an
additional degree shift. Thus it vanishes entirely if $b>0$;
if $b=0$ its least possible degree above $kd$ is
\[
 l(2p-3)=2k-3l\geq2k-3C.
\]
The case $k=0$ of Equation~\eqref{conn:cofiber-bound} follows
directly from $T_{p,0}(d)=S^0$.

For completeness, these mod-$p$ calculations imply the stated
homotopy bounds, not just cohomology bounds. Classifying spaces
of the finite groups used above have finitely many cells in
each degree, and the virtual Thom spectra are bounded below
after suspension. Their $p$-local integral homology groups,
and those of the indicated cofibers, are finitely generated in
each degree. Cohomology vanishing over the field $\F_p$ gives
homology vanishing in the same range. The universal coefficient
sequence then shows that each integral homology group in that
range has zero quotient modulo $p$. Finite generation over
$\Z_{(p)}$ forces that group to vanish. Finally, for a
bounded-below spectrum, the first nonzero homotopy group, if
there is one below the claimed bound, maps isomorphically to
the first nonzero integral homology group by the stable
Hurewicz theorem. This would contradict the integral homology
vanishing. The same argument in every degree proves
contractibility in the cases where all mod-$p$ cohomology
vanishes.
\end{proof}

\subsection{The bound for the universal Koszul module}

The estimates in Proposition~\ref{conn:symmetric} now apply
factor by factor to Equation~\eqref{eq:whole-module}. The crucial
point is that the number of variable types is fixed, although
their multiplicities may grow with the prime.

\begin{proposition}\label{conn:estimate}
Fix $0\leq a<n$ and $\lambda,\delta\in\Lambda$,
and set
\begin{equation}\label{conn:constants}
 C=\max\{0,\sigma(\lambda)\}+|\sigma(\delta)|+1,
 \qquad N_a=2^{a+1}-1,
 \qquad E=N_a(5C+1).
\end{equation}
On a set of primes in $\U$, and for all its sufficiently large
members, the underlying unfiltered module satisfies
\begin{equation}\label{conn:module-bound}
 K_p^a(p\lambda+\delta)
 \in\Sp_{\geq d(p\lambda+\delta)
                  +2\nu(p\lambda+\delta)-E}.
\end{equation}
The constants are independent of the individual multiplicities
occurring at this weight and of the prime.
\end{proposition}

\begin{proof}
Write $\gamma=p\lambda+\delta$. Restrict to the set of primes
where Equation~\eqref{eq:whole-module} holds for whole graded
modules. If $\gamma$ is outside the support monoid of the
universal model, its component is zero and there is nothing to
prove. Otherwise decompose the free algebra $\gr R_p^a$ as the
tensor product of the free algebras on the individual generator
spheres. There are
$\sum_{j=0}^a2^j=N_a$ such generators.

For a nonempty subset $S\subseteq[j)$, set $r=|S|$ as before.
The factor with multiplicity $k=k_{j,S}$ is
\[
 \Sym^k\big(S^{d(\gamma_{j,S})-r}\big).
\]
For $S=\varnothing$, the corresponding quotient factor in
Equation~\eqref{eq:whole-module} has multiplicity-$k$ component
$T_{p,k}(d_j)$. Here multiplicity includes the generator in
the source of the multiplication map, or equivalently the
possible Koszul cell, as well as the ordinary polynomial
copies. Its filtration order is $k$ on both sides of that
map. Distinct factors tensor together, and the weight-$\gamma$
component is the finite sum of the terms satisfying
\begin{equation}\label{conn:weight-identity}
 \sum_{j,S}k_{j,S}\gamma_{j,S}=\gamma.
\end{equation}

Every generator weight has positive integral mass. Taking mass
in Equation~\eqref{conn:weight-identity} therefore gives
\[
 0\leq k_{j,S}\leq\sigma(\gamma)
     =p\sigma(\lambda)+\sigma(\delta)\leq Cp.
\]
This also proves finiteness of the sum at each prime. For
$p>C$, Proposition~\ref{conn:symmetric} applies to every
factor. Since $d(\gamma_{j,S})$ is even, the odd-degree
improvement applies exactly when $r$ is odd. For $r=0$, the
multiplication-cofiber improvement applies instead.
Tensoring lower degree bounds adds them: after desuspending
each factor to a connective spectrum, their tensor product is
connective. The bound for each tensor term, relative to
$d(\gamma)$, is consequently at least
\begin{equation}\label{conn:sum-bound}
 -\sum_{j,S}r k_{j,S}
 +\sum_{r\text{ odd}} k_{j,S}
 +2\sum_{r=0}k_{j,S}-E.
\end{equation}
Indeed the odd factors lose at most $2C+1$ each, the
multiplication cofibers lose at most $3C$ each, and there are
at most $N_a$ factors; the displayed value of $E$ bounds their
total loss, including factors of multiplicity zero.

The elementary inequality
\[
 -r+\mathbf 1_{\{r\text{ odd}\}}
       +2\mathbf 1_{\{r=0\}}\geq2(1-r)
 \qquad(r\geq0)
\]
and the identity $\nu(\gamma_{j,S})=1-r$ turn
Equation~\eqref{conn:sum-bound} into the lower bound
$2\nu(\gamma)-E$. Thus every associated graded piece at
weight $\gamma$ lies in the range asserted in
Equation~\eqref{conn:module-bound}.

At this weight the filtration is zero above
$\sigma(\gamma)$ and constant in filtration degrees at most
zero. Starting at its zero top term, the cofiber sequences
relating successive filtration terms and associated graded
pieces prove the same bound for the unfiltered component.
This uses only closure of $\Sp_{\geq L}$ under extensions.
The filtration length is finite at each prime; no uniform
bound on that length is required.
\end{proof}

\subsection{Vanishing after the Koszul-module test}

Recall that $i=(\lambda,\delta,m)\in I$, that $q=2(p-1)$,
and that the layer comparison uses
\[
 d(Fi)=d(p\lambda+\delta)-qm.
\]
The following consequence retains the strict inequality in
$m$ that will determine the normalization range.

\begin{lemma}\label{conn:degree-gap}
For every fixed $i=(\lambda,\delta,m)$ with
$m>-\nu(\lambda)$, one has
\[
 Q(K_p^a)_i=0,\qquad \Phi(K_p^a)_i=0.
\]
Here $Q$ and $\Phi$ are applied to the primewise underlying
graded modules, as allowed by the arbitrary-object part of
the layer comparison.
\end{lemma}

\begin{proof}
Let $s\in\Z$ be any fixed homotopy degree. The $s$th homotopy
group of $Q(K_p^a)_i$ is computed by maps from the sphere of
ordinary degree $d(Fi)+s$ to $K_p^a(p\lambda+\delta)$.
Subtracting this requested degree from the lower bound of
Proposition~\ref{conn:estimate} gives
\begin{align*}
 &d(p\lambda+\delta)+2\nu(p\lambda+\delta)-E
       -\big(d(Fi)+s\big)\\
 &\hspace{1cm}=
 2p\big(\nu(\lambda)+m\big)
       +2\nu(\delta)-2m-E-s.
\end{align*}
Because $\nu(\lambda)+m$ is a positive integer, this expression
tends to $+\infty$ with $p$. The requested primewise homotopy
group therefore vanishes on a set in $\U$. It follows that
$\pi_s Q(K_p^a)_i=0$. This argument applies separately to
every fixed $s$, so the mapping spectrum is zero. The
arbitrary-object base-change comparison of
Equation~\eqref{eq:layer} identifies
\[
 \Phi(K_p^a)_i\simeq
 Q(K_p^a)_i\otimes_D(D//\varpi),
\]
which proves the second assertion.
\end{proof}

\begin{proposition}\label{conn:k-test}
For every $\lambda\in\Lambda$ and every integer
$m>-\nu(\lambda)$, there is a vanishing of
$\delta$-graded modules
\begin{equation}\label{eq:k-test}
 K^a\otimes_{R^a}
       \big(B(R_p^a)/t\big)_{\lambda,*,m}=0.
\end{equation}
The tensor product uses the $R^a$-action in the
$\delta$-grading. The notation $/t$ denotes the underlying
module cofiber of multiplication by $t$.
\end{proposition}

\begin{proof}
Fix $\lambda$ and $m$. Trivial-action maps in the
$\delta$-direction, followed by the primewise module action,
give the balanced pairing
\begin{equation}\label{conn:module-pairing}
 K^a\otimes_{R^a}\Phi(R_p^a)_{\lambda,*,m}
 \longrightarrow \Phi(K_p^a)_{\lambda,*,m}.
\end{equation}
In constructing it we use the chosen comparisons
$h_\Lambda(R_p^a)\simeq R^a$ and
$h_\Lambda(K_p^a)\simeq K^a$ from
Section~\ref{sec:universal}. The compatibility of the module
action with these comparisons makes the pairing balanced over
the indicated $\delta$-scalar action.

We verify that Equation~\eqref{conn:module-pairing} is an
equivalence using the finite unfiltered module cells of $K_p^a$.
For a free cell generated in fixed weight $\alpha$ and ordinary
sphere degree $d(\alpha)+u$, with $u\in\Z$ fixed, the pairing
identifies both sides with the shift of
$\Phi(R_p^a)_{\lambda,*,m}$ by weight $\alpha$ and chain degree
$u$. Concretely, the primewise module input at $\delta$ on the
right is
$S^{d(\alpha)+u}\otimes R_p^a(p\lambda+\delta-\alpha)$.
After applying $J^\vee\otimes\Tr$, the Picard grading moves
the trivially acting $S(\alpha)$ into the source, changing its index to
$(\lambda,\delta-\alpha,m)$ and leaving the fixed suspension
$u$. This is exactly the free-cell shift on the left.
Both sides of Equation~\eqref{conn:module-pairing} are exact in
the module variable: tensor products preserve finite cofibers,
and the mapping constructions are exact componentwise.
Induction over the finitely many cells now proves the
equivalence for $K_p^a$.

If $m>-\nu(\lambda)$, Lemma~\ref{conn:degree-gap} shows that
the right side is zero at every fixed $\delta$. Finally,
Equation~\eqref{eq:layer} identifies
$B(R_p^a)//t$ with $\Phi(R_p^a)$ compatibly with the scalar
action. Since $t$ has chain degree zero, its algebra quotient
has the multiplication cofiber $B(R_p^a)/t$ as underlying
module, with the identification specified by its universal
null. Substituting this comparison into
Equation~\eqref{conn:module-pairing} proves
Equation~\eqref{eq:k-test}.
\end{proof}

The bounds in this section are asserted separately at each
fixed lattice index and each fixed integer suspension. Their
prime thresholds may depend on those indices. This is enough
to prove the vanishing of the corresponding mapping objects;
it does not require intersecting infinitely many large sets
of primes. In contrast, the use of
Equation~\eqref{eq:whole-module} at the moving weight
$p\lambda+\delta$ was justified by its equivalence of whole
primewise modules on one large set.
